
\clearpage
\VerdierPhysicalPageAnchor{0124}%
\begingroup
\small
\setlength{\emergencystretch}{1.5em}%
\setlength{\abovedisplayskip}{4pt plus 2pt minus 2pt}%
\setlength{\belowdisplayskip}{4pt plus 2pt minus 2pt}%

\VerdierSemanticAnchor{ch2-1-3-2-triangle-morphism-conditions}
\noindent 使得 $(\operatorname{id}_{X_1},u_1,m'_1)$ 与
$(u_3,\operatorname{id}_{X_3},m'_3)$ 都是三角的态射，并且三角

\VerdierSemanticAnchor{ch2-1-3-2-connecting-triangle}
\[
  Z'_3\xrightarrow{m'_1}Z'_2\xrightarrow{m'_3}Z'_1
  \xrightarrow{\,v'_3[1]w'_1\,}Z'_3[1]
\]
\noindent 是特异三角。

\VerdierSemanticAnchor{ch2-1-3-2-comparison-isomorphisms}
此外，由 (1.2.4) 与 (1.2.14)，存在同构
$i_j\colon Z'_j\longrightarrow Z_j$（$1\leq j\leq 3$），使得
$(\operatorname{id},\operatorname{id},i_j)$ 是三角的同构。于是只需令

\VerdierSemanticAnchor{ch2-1-3-2-transported-morphisms}
\[
  m_1=i_2m'_1i_3^{-1}\colon Z_3\longrightarrow Z_2,
  \qquad
  m_3=i_1m'_3i_2^{-1}\colon Z_2\longrightarrow Z_1.
\]

\VerdierSemanticAnchor{ch2-1-3-2-proof-complete}
态射 $m_1$ 与 $m_3$ 具有公理 TRIV 所要求的性质。

\VerdierSemanticAnchor{ch2-1-3-3}
\medskip
\noindent\textbf{1.3.3.}\quad
设 $C$ 为一个\emph{加法}范畴。对 $C$ 的每个对象 $X$，关联 $C$ 中如下复形：
除零次分量等于 $X$ 外，其余分量均为零。由此定义一个从 $C$ 到
$\operatorname{comp}(C)$ 的函子；再与典范函子
$\operatorname{comp}(C)\longrightarrow K(C)$ 复合，得到从 $C$ 到 $K(C)$
的函子。

\VerdierSemanticAnchor{ch2-1-3-3-full-embeddings}
这些函子都是全忠实的，并且在对象上单射，因而把 $C$ 实现为
$\operatorname{comp}(C)$ 与 $K(C)$ 的满子范畴。此后，总是通过这些函子把
范畴 $C$ 看作 $\operatorname{comp}(C)$ 与 $K(C)$ 的满子范畴。

\VerdierSemanticAnchor{ch2-1-3-4}
\medskip
\noindent\textbf{1.3.4.}\quad
设 $C$ 为一个\emph{阿贝尔}范畴。对于 $C$ 中每个复形 $X$，用 $H^0(X)$
表示 $C$ 中的对象

\VerdierSemanticAnchor{ch2-1-3-4-1}
\[
  H^0(X)=\operatorname{Ker}(d_X^0)/\operatorname{Im}(d_X^{-1}).
  \tag{1.3.4.1}
\]

\VerdierSemanticAnchor{ch2-1-3-4-induced-morphism}
对于每个复形态射 $X\xrightarrow{f}Y$，用
\[
  H^0(f)\colon H^0(X)\longrightarrow H^0(Y)
\]
\noindent 表示由 $f$ 诱导的态射。由此定义一个函子

\VerdierSemanticAnchor{ch2-1-3-4-comp-functor}
\[
  H^0\colon\operatorname{comp}(C)\longrightarrow C
\]
\noindent 称为\emph{第零上同调对象}。该函子把同伦态射送到相等的态射，
从而诱导出一个仍记为

\VerdierSemanticAnchor{ch2-1-3-4-k-functor}
\[
  H^0\colon K(C)\longrightarrow C.
\]
\endgroup

